\documentclass[11pt]{article}

%%%%%%%% PAQUETES %%%%%%%%
\usepackage[spanish, es-tabla]{babel}       % 'Tabla' en vez de 'Cuadro'
\usepackage[T1]{fontenc}
\usepackage{lmodern}                               % Texto con serifa
\usepackage{tgheros}                               % Títulos sin serifa (Helvetica)
\usepackage{microtype}
\usepackage[letterpaper, margin=2.5cm]{geometry}   % Tamaño carta
\usepackage{graphicx}
\usepackage[table]{xcolor}
\usepackage{tikz}
\usetikzlibrary{calc, shapes.symbols, positioning}
\usepackage{tocloft}                               % Estilo del índice
\usepackage{titlesec}                              % Estilo de títulos
\usepackage{fancyhdr}                              % Encabezado y pie
\usepackage{tabularx}                              % Tablas a todo el ancho
\usepackage{booktabs}                              % Líneas de tabla elegantes
\usepackage{array}
\usepackage{enumitem}                              % Listas compactas
\usepackage{float}                                 % Figuras con [H]
\usepackage{caption}
\usepackage{subcaption}                            % Figuras lado a lado
\usepackage{listings}                              % Bloques de comandos
\usepackage{etoolbox}
\usepackage{needspace}                             % Evita títulos solos al final de página
\usepackage[most]{tcolorbox}                       % Recuadros de "pendiente"
\usepackage{hyperref}

%%%%%%%% COLORES %%%%%%%%
\definecolor{ucbazul}{HTML}{0B2A5B}
\definecolor{ucbazulclaro}{HTML}{1F4E8C}
\definecolor{acento}{HTML}{C9A227}
\definecolor{grisTexto}{HTML}{4A4A4A}
\definecolor{fondoCodigo}{HTML}{F4F4F4}
\definecolor{verdeOK}{HTML}{2E7D32}

\hypersetup{colorlinks=true, linkcolor=ucbazul, linktoc=page, urlcolor=ucbazulclaro,
            pdftitle={Proyecto Final - Performance Testing con k6}}

%%%%%%%% DATOS DEL PROYECTO (editar aquí si cambian) %%%%%%%%
\newcommand{\repoURL}{https://github.com/Soliz-G-David/quickpizza-performance-k6}
\newcommand{\appURL}{https://quickpizza.grafana.com}

% Texto en mayúsculas con espaciado entre letras
\newcommand{\espaciado}[1]{\textls[120]{\MakeUppercase{#1}}}

% Captura con marco; si la imagen aún no existe no se muestra nada
% Uso: \captura{images/archivo.png}{Pie de figura}
\newcommand{\captura}[2]{%
  \IfFileExists{#1}{%
  \begin{figure}[!htbp]
    \centering
    \fcolorbox{gray!45}{white}{\includegraphics[width=\dimexpr\textwidth-2\fboxsep-2\fboxrule\relax, height=0.40\textheight, keepaspectratio]{#1}}%
    \caption{#2}
  \end{figure}}{}}

% Igual que \captura pero con etiqueta para referenciarla
\newcommand{\capturaref}[3]{%
  \IfFileExists{#1}{%
  \begin{figure}[!htbp]
    \centering
    \fcolorbox{gray!45}{white}{\includegraphics[width=\dimexpr\textwidth-2\fboxsep-2\fboxrule\relax, height=0.40\textheight, keepaspectratio]{#1}}%
    \caption{#2}
    \label{#3}
  \end{figure}}{}}

% Recuadro dorado para lo que falta completar
\newtcolorbox{pendiente}{enhanced, colback=acento!8, colframe=acento, boxrule=0pt, leftrule=3pt,
  arc=0pt, left=8pt, right=8pt, top=5pt, bottom=5pt, fontupper=\small\sffamily,
  before skip=1.3em,
  overlay={\node[anchor=south east, fill=acento, text=white, font=\sffamily\bfseries\tiny,
    inner sep=2.5pt] at (frame.north east) {\espaciado{Pendiente}};}}

%%%%%%%% ESTILO DE FIGURAS, TABLAS Y CÓDIGO %%%%%%%%
\captionsetup{font=small, labelfont={bf,sf,color=ucbazul}, labelsep=period, justification=centering}
\setlist{nosep, leftmargin=*}
\setlist[itemize]{label={\color{acento}\rule[0.2ex]{0.45em}{0.45em}}}
\renewcommand{\arraystretch}{1.35}
\setlength{\aboverulesep}{0pt}                      % Sin huecos blancos entre filas coloreadas
\setlength{\belowrulesep}{0pt}
\AtBeginEnvironment{tabularx}{\rowcolors{2}{white}{ucbazul!4}}  % Filas alternas suaves
\AtBeginEnvironment{table}{\small}                  % Tablas en letra un poco menor

% Insignias de estado: \aprobado y \noaprobado
\newcommand{\insignia}[3]{\tcbox[on line, colback=#1, colframe=#1, arc=6pt, boxsep=0pt,
  left=5pt, right=5pt, top=1.5pt, bottom=1.5pt]{\sffamily\scriptsize\bfseries\color{#2}#3}}
\newcommand{\aprobado}{\insignia{verdeOK!12}{verdeOK}{APROBADO}}
\newcommand{\noaprobado}[1]{\insignia{red!10}{red!70!black}{NO APROBADO}\ {\scriptsize #1}}
% Etiquetas de color por escenario
\newcommand{\chip}[2]{\tcbox[on line, colback=#1, colframe=#1, arc=2pt, boxsep=0pt,
  left=4pt, right=4pt, top=1.5pt, bottom=1.5pt]{\sffamily\scriptsize\bfseries\color{white}#2}}
\newcommand{\esmoke}{\chip{gray!70}{SMOKE}}
\newcommand{\eload}{\chip{ucbazulclaro}{LOAD}}
\newcommand{\espike}{\chip{acento!90!black}{SPIKE}}
\newcommand{\estress}{\chip{red!65!black}{STRESS}}
% Subtítulo con letra en un círculo dorado. Uso: \bloque{A}{Título}
\newcommand{\bloque}[2]{\par\needspace{6\baselineskip}\vspace{0.8em}\noindent
  \tikz[baseline=(l.base)]\node[circle, fill=acento, text=white, font=\sffamily\bfseries\small,
    inner sep=1.5pt, minimum size=1.35em](l){#1};%
  \hspace{0.5em}{\sffamily\bfseries\normalsize\color{ucbazulclaro}#2}\par\vspace{0.3em}}

% Tablas estilo B: encabezado azul con texto blanco y línea dorada al final
\newcommand{\encab}[1]{{\leavevmode\color{white}\sffamily\bfseries #1}}
\renewcommand{\toprule}{\arrayrulecolor{ucbazul}\specialrule{0.8pt}{0pt}{0pt}\arrayrulecolor{black}}
\renewcommand{\midrule}{\arrayrulecolor{ucbazul}\specialrule{0.4pt}{0pt}{0pt}\arrayrulecolor{black}}
\renewcommand{\bottomrule}{\arrayrulecolor{acento}\specialrule{1.2pt}{0pt}{0pt}\arrayrulecolor{black}}

%%%%%%%% ELEMENTOS NUEVOS %%%%%%%%
% Etiqueta pequeña en mayúsculas doradas
\newcommand{\etiqueta}[1]{{\sffamily\bfseries\footnotesize\color{acento}\espaciado{#1}}}

% Resumen ejecutivo (recuadro azul)
\newtcolorbox{resumen}{enhanced, colback=ucbazul, colframe=ucbazul, arc=3pt, coltext=white,
  left=12pt, right=12pt, top=8pt, bottom=8pt, fontupper=\sffamily\small,
  title={\etiqueta{Resumen ejecutivo}}, attach title to upper, after title={\par\vspace{3pt}}}

% Tarjeta con un indicador clave. Uso: \kpi{Nombre}{Valor}{Detalle}
\newcommand{\kpi}[3]{%
  \begin{tcolorbox}[enhanced, width=0.235\textwidth, nobeforeafter, colback=white,
    colframe=ucbazul!15, boxrule=0.6pt, arc=3pt, borderline north={3pt}{0pt}{acento},
    height=2.05cm, valign=center,
    halign=center, top=6pt, bottom=6pt]
  {\sffamily\scriptsize\color{grisTexto}\espaciado{#1}}\\[2pt]
  {\sffamily\bfseries\fontsize{18}{22}\selectfont\color{ucbazul}#2}\\[1pt]
  {\sffamily\scriptsize\color{verdeOK}#3}
  \end{tcolorbox}}

% Recuadro para hallazgos importantes
\newtcolorbox{hallazgo}{enhanced, colback=ucbazulclaro!6, colframe=ucbazulclaro!6, arc=0pt,
  borderline west={4pt}{0pt}{ucbazulclaro}, left=10pt, right=8pt, top=6pt, bottom=6pt,
  fontupper=\small}

% Gráfico del perfil de carga. Uso: \perfil{Título}{VUs máx.}{segundos}{coordenadas (s,VUs)}{marcas eje x}
\newcommand{\perfil}[5]{%
  \begin{tikzpicture}[x=4.2cm/#3, y=2.2cm/#2]
    \fill[ucbazul!8] (0,0) -- plot coordinates {#4} -- (#3,0) -- cycle;
    \draw[ucbazul, line width=1.4pt, line join=round] plot[mark=*, mark size=1.4pt, mark options={fill=acento, draw=acento}] coordinates {#4};
    \draw[gray!60] (0,0) -- (#3*1.04,0);
    \draw[gray!60] (0,0) -- (0,#2*1.12);
    \draw[gray!30, dashed] (0,#2) -- (#3,#2);
    \node[left, font=\sffamily\scriptsize, text=grisTexto] at (0,#2) {#2};
    \node[left, font=\sffamily\scriptsize, text=grisTexto] at (0,0) {0};
    \foreach \t/\l in {#5} \node[below, font=\sffamily\scriptsize, text=grisTexto] at (\t,0) {\l};
    \node[font=\sffamily\bfseries\small, text=ucbazul] at (#3/2,#2*1.3) {#1};
  \end{tikzpicture}}

\lstset{
  basicstyle=\ttfamily\small,
  backgroundcolor=\color{fondoCodigo},
  frame=leftline, framerule=2pt, rulecolor=\color{acento},
  xleftmargin=0.4cm, framexleftmargin=0.3cm,
  breaklines=true, columns=fullflexible, keepspaces=true,
  aboveskip=0.6em, belowskip=0.8em, showstringspaces=false,
  % Tildes y eñes dentro del código
  literate={á}{{\'a}}1 {é}{{\'e}}1 {í}{{\'i}}1 {ó}{{\'o}}1 {ú}{{\'u}}1
           {Á}{{\'A}}1 {É}{{\'E}}1 {Í}{{\'I}}1 {Ó}{{\'O}}1 {Ú}{{\'U}}1
           {ñ}{{\~n}}1 {Ñ}{{\~N}}1 {ü}{{\"u}}1
}

%%%%%%%% CONSOLA (GIT BASH) Y ARCHIVOS DE CÓDIGO %%%%%%%%
\definecolor{bashFondo}{HTML}{0C0C0C}
\definecolor{bashTexto}{HTML}{CCCCCC}
\definecolor{bashVerde}{HTML}{16C60C}
\definecolor{bashMorado}{HTML}{B148C6}
\definecolor{bashAmarillo}{HTML}{C19C00}
\definecolor{bashComentario}{HTML}{767676}

% Línea del prompt de Git Bash
\newcommand{\prompt}{{\color{bashVerde}PC- DAVID@LAPTOP-QNOIOQBU} {\color{bashMorado}MINGW64} {\color{bashAmarillo}\textasciitilde/quickpizza-performance-k6}}

% Icono de Git (rombo naranja) y botones de la ventana
\newcommand{\iconoGit}{\tikz[baseline=-0.6ex]{\fill[orange!85!red, rotate=45] (-0.12,-0.12) rectangle (0.12,0.12);
  \draw[white, line width=0.6pt] (-0.06,-0.05) -- (0.02,0.03) (0.02,0.03) -- (0.07,-0.02);}}
\newcommand{\botonesVentana}{{\sffamily\footnotesize\color{black!70}\rule[0.35ex]{0.6em}{0.5pt}\hspace{1.4em}%
  \tikz[baseline=-0.1ex]\draw[black!70, line width=0.5pt] (0,0) rectangle (0.2,0.2);\hspace{1.3em}$\times$}}

% Consola Git Bash. Uso: \begin{gitbash} ... \end{gitbash}
% Cada comando empieza con "$ ". El prompt se agrega solo al inicio.
\newtcblisting{gitbash}{listing only, enhanced, colback=bashFondo, colframe=black!35,
  boxrule=0.5pt, arc=0pt, left=8pt, right=8pt, top=3pt, bottom=4pt,
  before skip=0.8em, after skip=1.2em,
  title={\iconoGit\hspace{0.6em}{\sffamily\footnotesize\color{black!85}MINGW64:/c/Users/PC- DAVID/quickpizza-performance-k6}\hfill\botonesVentana},
  colbacktitle=white!96!black, coltitle=black, toptitle=3pt, bottomtitle=3pt, titlerule=0pt,
  listing options={basicstyle=\ttfamily\small\color{bashTexto}, backgroundcolor={}, frame=none,
    xleftmargin=0pt, framexleftmargin=0pt, breaklines=true, columns=fullflexible, keepspaces=true,
    escapeinside={(*}{*)}, upquote=true,
    morecomment={[l]{\#}}, commentstyle=\color{bashComentario}\itshape,
    aboveskip=0pt, belowskip=0pt}}

% Archivo de código con barra de título. Uso: \begin{archivo}{utils/data.js} ... \end{archivo}
\lstdefinelanguage{JavaScript}{
  keywords={export, import, from, const, let, var, new, function, return, if, else, true, false},
  morekeywords=[2]{SharedArray, open, Number, parseCsv},
  sensitive=true, morecomment=[l]{//}, morestring=[b]"}
\newtcblisting{archivo}[2][]{listing only, enhanced, breakable, colback=white, colframe=ucbazul,
  arc=0pt, boxrule=0pt, leftrule=3pt, left=8pt, right=6pt, top=3pt, bottom=3pt,
  borderline={0.4pt}{0pt}{ucbazul!20}, before skip=0.7em, after skip=0.9em,
  title={\sffamily\scriptsize\bfseries\espaciado{#2}}, colbacktitle=ucbazul, coltitle=white,
  toptitle=2pt, bottomtitle=2pt,
  listing options={basicstyle=\ttfamily\small, keywordstyle=\color{ucbazulclaro}\bfseries,
    keywordstyle={[2]\color{ucbazul}}, stringstyle=\color{verdeOK}, commentstyle=\color{gray}\itshape,
    backgroundcolor={}, frame=none, xleftmargin=0pt, framexleftmargin=0pt, breaklines=true,
    columns=fullflexible, keepspaces=true, upquote=true, aboveskip=0pt, belowskip=0pt, #1}}

%%%%%%%% ESTILO DEL ÍNDICE %%%%%%%%
\renewcommand{\cfttoctitlefont}{\sffamily\bfseries\fontsize{28}{32}\selectfont\color{ucbazul}}
\renewcommand{\cftaftertoctitle}{\par\vspace{0.2cm}{\color{acento}\rule{\linewidth}{1.4pt}}}
\setlength{\cftbeforetoctitleskip}{0pt}
\setlength{\cftaftertoctitleskip}{0.6cm}

% Secciones
\renewcommand{\cftsecfont}{\sffamily\bfseries\large\color{ucbazul}}
\renewcommand{\cftsecpagefont}{\sffamily\bfseries\large\color{ucbazul}}
\renewcommand{\cftsecleader}{\hfill}
\setlength{\cftbeforesecskip}{0.55em}
% Número de sección en un cuadro azul
\newsavebox{\cajaindice}
\renewcommand{\cftsecpresnum}{\begin{lrbox}{\cajaindice}\color{white}}
\renewcommand{\cftsecaftersnum}{\end{lrbox}\colorbox{ucbazul}{\makebox[1em]{\usebox{\cajaindice}}}}
\setlength{\cftsecnumwidth}{2.2em}

% Subsecciones
\renewcommand{\cftsubsecfont}{\small\color{grisTexto}}
\renewcommand{\cftsubsecpagefont}{\small\color{grisTexto}}
\renewcommand{\cftsubsecpresnum}{\color{acento!85!black}}
\renewcommand{\cftsubsecleader}{\color{gray!60}\cftdotfill{\cftdotsep}}
\setlength{\cftbeforesubsecskip}{0.05em}
\setlength{\cftsubsecindent}{2.2em}
\setlength{\cftsubsecnumwidth}{2.6em}

%%%%%%%% ESTILO DE TÍTULOS %%%%%%%%
% Sección: número en un cuadro azul y línea dorada debajo de todo el ancho
\titleformat{\section}[block]
  {\sffamily\bfseries\Large\color{ucbazul}}
  {\colorbox{ucbazul}{\makebox[1.1em]{\color{white}\thesection}}}{0.6em}{}
  [\vspace{0.3em}{\color{acento}\titlerule[1.2pt]}]
\titleformat{\subsection}
  {\sffamily\bfseries\large\color{ucbazulclaro}}
  {\thesubsection}{0.6em}{}
\titleformat{\subsubsection}
  {\sffamily\bfseries\normalsize\color{ucbazulclaro}}{}{0pt}{}
\titlespacing*{\section}{0pt}{1.6em}{1em}
\titlespacing*{\subsection}{0pt}{1.2em}{0.6em}

%%%%%%%% ENCABEZADO Y PIE %%%%%%%%
\usepackage{lastpage}                              % Total de páginas
\pagestyle{fancy}
\fancyhf{}
\fancyhead[L]{\IfFileExists{images/images.png}{\raisebox{-0.25\height}{\includegraphics[height=0.75cm]{images/images.png}}}{}}
\fancyhead[R]{\sffamily\small\color{grisTexto}Performance Testing con k6 -- Grupo 2}
\fancyfoot[L]{\sffamily\scriptsize\color{grisTexto}Diplomado en Software Testing -- UCB}
\fancyfoot[R]{\sffamily\scriptsize\color{grisTexto}Página {\color{ucbazul}\bfseries\thepage} de \pageref*{LastPage}}
\renewcommand{\headrulewidth}{0.8pt}
\renewcommand{\footrulewidth}{0.4pt}
\renewcommand{\footrule}{{\color{acento!60}\hrule width\headwidth height\footrulewidth \vskip-\footrulewidth}}
\renewcommand{\headrule}{{\color{acento}\hrule width\headwidth height\headrulewidth \vskip-\headrulewidth}}
\setlength{\headheight}{30pt}
\fancypagestyle{plain}{}   % Las páginas 'plain' (índice) usan el mismo estilo

% Las figuras flotan para llenar la página y evitar espacios en blanco
\renewcommand{\topfraction}{0.9}
\renewcommand{\bottomfraction}{0.9}
\renewcommand{\textfraction}{0.07}
\renewcommand{\floatpagefraction}{0.75}
\setcounter{topnumber}{3}
\setcounter{bottomnumber}{3}
\setcounter{totalnumber}{5}
\makeatletter
\setlength{\@fptop}{0pt}          % Páginas de solo figuras: empiezan arriba, sin centrar
\setlength{\@fpsep}{1.2em}
\makeatother

\setlength{\parskip}{0.5em}
\setlength{\parindent}{0pt}

%%%%%%%% DOCUMENTO %%%%%%%%
\begin{document}

\pagenumbering{roman}
%%%% HOJA 1: CARÁTULA
\begin{titlepage}
\thispagestyle{empty}

% --- Fondo: banda inferior y marca de agua ---
\begin{tikzpicture}[remember picture, overlay]
  \fill[ucbazul] (current page.south west) rectangle ($(current page.south east)+(0,1.1cm)$);
  \fill[acento]  ($(current page.south west)+(0,1.1cm)$) rectangle ($(current page.south east)+(0,1.2cm)$);
  \node[text=white] at ($(current page.south)+(0,0.55cm)$)
    {\sffamily\footnotesize \espaciado{Departamento de Ingeniería de Sistemas \quad·\quad Cochabamba -- Bolivia}};
  \node[anchor=east, text=ucbazul!5] at ($(current page.east)+(0.6cm,-2.2cm)$)
    {\sffamily\bfseries\fontsize{260}{260}\selectfont k6};
\end{tikzpicture}

% --- Encabezado: logo a la izquierda y programa a la derecha ---
\vspace*{0.1cm}
\noindent
\begin{minipage}[c]{0.47\linewidth}
\IfFileExists{images/images.png}{%
  \includegraphics[height=3.3cm, keepaspectratio]{images/images.png}%
}{%
  \tikz\draw[gray!60, dashed, rounded corners] (0,0) rectangle (7cm,3.3cm)
    node[midway, text=gray, font=\small] {Logo UCB};%
}
\end{minipage}%
\hfill
\begin{minipage}[c]{0.49\linewidth}
\begin{tcolorbox}[enhanced, frame hidden, colback=white, arc=0pt, borderline west={1.5pt}{0pt}{acento},
  left=16pt, right=0pt, top=4pt, bottom=4pt, boxsep=0pt]
{\sffamily\small\color{grisTexto}\espaciado{Diplomado en}}\\[0.1cm]
{\sffamily\bfseries\fontsize{24}{28}\selectfont\color{ucbazul}Software Testing}\\[0.3cm]
\begin{tikzpicture}[baseline=(n.base)]
  \node[anchor=base west, font=\sffamily\bfseries\fontsize{46}{46}\selectfont, text=acento, inner sep=0pt] (n) {05};
  \node[anchor=south west, font=\sffamily\bfseries\scriptsize, text=grisTexto, inner sep=0pt]
    at ($(n.north west)+(0.05cm,0.1cm)$) {\espaciado{Módulo}};
  \draw[ucbazul, line width=1.2pt] ($(n.north east)+(0.3cm,0.05cm)$) -- ($(n.south east)+(0.3cm,0)$);
  \node[anchor=west, align=left, inner sep=0pt] at ($(n.east)+(0.6cm,0)$)
    {{\sffamily\bfseries\Large\color{ucbazulclaro}Performance}\\[0.05cm]
     {\sffamily\bfseries\Large\color{ucbazulclaro}Testing}};
\end{tikzpicture}
\end{tcolorbox}
\end{minipage}

\vspace{0.5cm}
{\color{acento}\rule{\linewidth}{1pt}}

% --- Título ---
\vspace{2.2cm}
{\sffamily\bfseries\huge\color{acento}\espaciado{Proyecto Final}}\\[0.3cm]
{\sffamily\bfseries\fontsize{38}{42}\selectfont\color{ucbazul} Performance Testing\\[0.1cm] con k6}\\[0.5cm]
{\sffamily\Large\color{grisTexto} Evaluación de rendimiento de QuickPizza\\ hasta 400 usuarios virtuales}

% --- Datos ---
\vfill
\begin{tcolorbox}[enhanced, colback=ucbazul!4, colframe=ucbazul!4, arc=2pt, left=12pt, right=12pt,
  top=10pt, bottom=10pt, borderline west={3pt}{0pt}{acento}]
\large
\renewcommand{\arraystretch}{1.25}
\begin{tabular}{@{}r@{\hspace{0.6cm}}l@{}}
  {\sffamily\bfseries\color{ucbazul}Integrantes} & Pablo Ivan Chuquimia Huanca \\
                                                & Mauricio Gutierrez Pacheco \\
                                                & Rodrigo Camilo Mendez Gallardo \\
                                                & Maria Alexandra Quintanilla \\
                                                & Pedro Adrian Santos Orellana \\
                                                & David Soliz Garcia \\[0.25cm]
  {\sffamily\bfseries\color{ucbazul}Grupo}       & 2 \\[0.25cm]
  {\sffamily\bfseries\color{ucbazul}Docente}     & Marco Luna Gonzales \\[0.25cm]
  {\sffamily\bfseries\color{ucbazul}Fecha}       & 16 de noviembre de 2026 \\
\end{tabular}
\end{tcolorbox}
\vspace{0.6cm}
\end{titlepage}

%%%% HOJA 2: ÍNDICE
\pagenumbering{roman}
\setcounter{page}{2}
\tableofcontents
\clearpage

%%%% CONTENIDO
\pagenumbering{arabic}

%%%% RESUMEN EJECUTIVO E INDICADORES CLAVE
\begin{resumen}
QuickPizza cumplió los criterios de éxito bajo carga nominal: con 100 usuarios
virtuales el P95 fue de 248.5\,ms (umbral de 500\,ms), con 0\,\% de errores y
100\,\% de checks, y absorbió un pico de 100 VUs en 10 segundos sin degradarse.
Con 400 VUs se mantuvo sin errores (P95 de 278\,ms), pero se acerca a su límite:
la capacidad recomendada ronda los 300 usuarios concurrentes.
\end{resumen}

\vspace{0.3cm}
\noindent
\kpi{P95 con 100 VUs}{248.5 ms}{umbral 500 ms}\hfill
\kpi{Errores}{0 \%}{en las 4 pruebas}\hfill
\kpi{Checks}{100 \%}{23\,282 / 23\,282 (load)}\hfill
\kpi{Pico de carga}{330 req/s}{con 400 VUs}
\vspace{0.4cm}

%==================================================================
\section{Objetivo del trabajo}
%==================================================================

El objetivo de este proyecto es \textbf{evaluar de forma integral el rendimiento,
la estabilidad y el comportamiento bajo carga} de la aplicación web
\textbf{QuickPizza} mediante scripts automatizados en \textbf{Grafana k6}.

Para lograrlo, el equipo siguió un proceso incremental:

\begin{enumerate}
  \item \textbf{Grabación} del recorrido real de un usuario con \textbf{Grafana k6 Studio}:
        ingreso a la página, inicio de sesión, pedido de pizzas y calificación
        (\emph{No thanks} y \emph{Love it}).
  \item \textbf{Refinamiento del script} generado: correlación de valores dinámicos
        (token CSRF, token de sesión e ID de la pizza) y parametrización de datos con
        \texttt{SharedArray} (credenciales y 30 pizzas distintas).
  \item \textbf{Estructuración modular} del código en un repositorio Git, separando
        datos, utilidades, configuración y escenarios.
  \item \textbf{Ejecución de escenarios de carga} (\emph{smoke}, carga sostenida, \emph{spike}
        y estrés) desde 1 hasta \textbf{400 usuarios virtuales}, con validaciones y umbrales
        de aprobación.
  \item \textbf{Observabilidad en tiempo real} con el \textbf{dashboard web de Grafana k6},
        ejecutado en entorno local, y exportación de reportes HTML interactivos.
  \item \textbf{Análisis de resultados}, identificación de cuellos de botella y
        recomendaciones de optimización.
\end{enumerate}

%==================================================================
\section{Introducción y alcance}
%==================================================================

\subsection{Contexto de la aplicación}

\textbf{QuickPizza} (\url{\appURL}) es la aplicación web de demostración que
Grafana Labs publica para practicar pruebas de rendimiento con k6. Permite
registrarse e iniciar sesión, pedir una recomendación de pizza al azar según
filtros (calorías, ingredientes, si es vegetariana) y calificar las pizzas
recomendadas.

\subsection{Arquitectura tecnológica}

QuickPizza tiene un \emph{frontend} web que consume una \textbf{API REST} con
respuestas JSON. La autenticación usa un \textbf{token CSRF} previo al login; el
login devuelve un \textbf{token de sesión} que viaja en una \emph{cookie} y en la
cabecera \texttt{Authorization}. Las recomendaciones y las calificaciones se manejan
con los recursos \texttt{/api/pizza} y \texttt{/api/ratings}.

El sitio público funciona con \textbf{varias instancias (réplicas)} detrás de un
balanceador de carga. Este detalle se observó durante las pruebas (ver
sección~\ref{sec:parametrizacion}).

\subsection{Flujos críticos de negocio}

\begin{itemize}
  \item \textbf{Autenticación:} obtención del \texttt{csrf\_token} e inicio de sesión.
  \item \textbf{Recomendación de pizza:} \texttt{POST /api/pizza} con los filtros del usuario.
  \item \textbf{Calificación:} \texttt{POST /api/ratings} con el \texttt{pizza\_id} recomendado (\emph{Love it} o \emph{No thanks}).
\end{itemize}

\subsection{Criterios de éxito y objetivos de rendimiento}

Se adoptaron los criterios de referencia del enunciado como objetivos de
rendimiento del proyecto. Se aplican por igual a todos los escenarios:

\begin{table}[H]
\centering
\caption{Criterios de éxito.}
\begin{tabularx}{\textwidth}{l l >{\raggedright\arraybackslash}X}
\toprule
\rowcolor{ucbazul}
\encab{Métrica} & \encab{Objetivo} & \encab{Significado} \\
\midrule
\texttt{http\_req\_duration} & p(95) < 500\,ms & El 95\,\% de las peticiones responde en menos de medio segundo \\
\texttt{http\_req\_failed}   & rate < 1\,\%     & Menos de 1 de cada 100 peticiones falla \\
\texttt{checks}              & 100\,\%          & Todas las validaciones de negocio se cumplen \\
\bottomrule
\end{tabularx}
\end{table}

Como referencia previa, la prueba funcional del script v3 (3 VUs, 6 iteraciones)
dio P95 = 234\,ms y 0\,\% de errores. Los cuatro escenarios V5 cumplieron los tres
criterios (ver sección~\ref{sec:resultados}).

%==================================================================
\section{Funcionalidades y modelo de carga}
%==================================================================

\subsection{Flujos de usuario seleccionados}

Cada iteración de un usuario virtual ejecuta la \textbf{transacción completa},
organizada en seis grupos:

\begin{figure}[H]
\centering
\begin{tikzpicture}[paso/.style={signal, signal to=east, signal from=west, fill=ucbazul,
    draw=white, line width=1.5pt, text=white, font=\sffamily\bfseries\scriptsize, align=center,
    minimum height=1.05cm, minimum width=2.45cm, inner xsep=3pt},
    califica/.style={paso, fill=ucbazulclaro!75}, node distance=-0.05cm]
  \node[paso, signal from=nowhere] (a) {01\\HomePage};
  \node[paso, right=of a] (b) {02\\Login};
  \node[paso, right=of b] (c) {03\\Pedir pizza 1};
  \node[califica, right=of c] (d) {04\\Calificar 1};
  \node[paso, right=of d] (e) {05\\Pedir pizza 2};
  \node[califica, right=of e, fill=verdeOK] (f) {06\\Calificar 2};
  \node[below=0.12cm of b, font=\sffamily\tiny, text=grisTexto] {csrf + token de sesión};
  \node[below=0.12cm of c, font=\sffamily\tiny, text=grisTexto] {pizza.id};
  \node[below=0.12cm of d, font=\sffamily\tiny, text=grisTexto] {Love it / No thanks};
  \node[below=0.12cm of f, font=\sffamily\tiny, text=grisTexto] {Love it / No thanks};
\end{tikzpicture}
\caption{Flujo de negocio de cada iteración.}
\end{figure}

\begin{table}[H]
\centering
\caption{Grupos del flujo de negocio y peticiones principales.}
\begin{tabularx}{\textwidth}{l l >{\raggedright\arraybackslash}X}
\toprule
\rowcolor{ucbazul}
\encab{Grupo} & \encab{Acción del usuario} & \encab{Peticiones principales} \\
\midrule
01 & HomePage       & \texttt{GET /}, \texttt{GET /api/config}, \texttt{GET /api/quotes}, \texttt{GET /api/tools} \\
02 & Login          & \texttt{GET /login}, \texttt{POST /api/csrf-token}, \texttt{POST /api/users/token/login}, \texttt{POST /api/users/token/authenticate} \\
03 & Pedir pizza 1      & \texttt{POST /api/pizza} \\
04 & Calificar pizza 1  & \texttt{POST /api/ratings} (\emph{Love it} o \emph{No thanks}, pizza 1) \\
05 & Pedir pizza 2      & \texttt{POST /api/pizza} \\
06 & Calificar pizza 2  & \texttt{POST /api/ratings} (\emph{Love it} o \emph{No thanks}, pizza 2) \\
\bottomrule
\end{tabularx}
\end{table}

Entre cada grupo se aplica un \emph{think time} aleatorio de 1 a 3 segundos para
simular a una persona real. Por eso cada iteración dura entre 10 y 20\,s
(unos 15\,s de media), aunque cada petición responde en unos 200\,ms.

Desde la v4, la calificación de las pizzas sigue una distribución: k6 cuenta los
resultados con las métricas personalizadas \texttt{love\_it} y \texttt{no\_thanks}.
Cada iteración registra dos calificaciones, y en las corridas se observó una
proporción aproximada de 80\,\% \emph{Love it} y 20\,\% \emph{No thanks}.

\subsection{Perfil de carga y tipos de prueba}

Los escenarios de la versión V5 reutilizan el mismo flujo de negocio
(\texttt{scenarios/\allowbreak orderPizza.js}). Cada script solo define sus propias
\texttt{options}: usuarios virtuales, etapas y umbrales.

\begin{table}[H]
\centering
\caption{Escenarios de prueba (V5).}
\begin{tabularx}{\textwidth}{l l c c >{\raggedright\arraybackslash}X}
\toprule
\rowcolor{ucbazul}
\encab{Escenario} & \encab{Script} & \encab{VUs máx.} & \encab{Duración} & \encab{Propósito} \\
\midrule
\esmoke  & \texttt{quickpizza-v5-smoke.js}  & 1   & 1 iteración          & Validar el flujo antes de cargar \\
\eload   & \texttt{quickpizza-v5-load.js}   & 100 & 5\,min (3 etapas)    & Carga nominal sostenida \\
\espike  & \texttt{quickpizza-v5-spike.js}  & 100 & 2\,min 10\,s (4 etapas) & Subida brusca de usuarios \\
\estress & \texttt{quickpizza-v5-stress.js} & 400 & 5\,min (5 etapas)    & Superar la carga nominal \\
\bottomrule
\end{tabularx}
\end{table}

El \emph{spike} sube de 9 a 100 VUs en solo 10 segundos (unos 9 usuarios nuevos por
segundo), los sostiene durante un minuto y luego baja. Su pico es igual al máximo
del \emph{load}: se diferencia por la velocidad de llegada de los usuarios, no por
la cantidad. El \emph{stress} sube en escalones de un minuto (100, 200, 300 y
400 VUs) y luego baja.

\begin{figure}[H]
\centering
\perfil{Load}{100}{300}{(0,0) (60,100) (242,100) (300,0)}{0/0, 60/1, 120/2, 180/3, 240/4, 300/5 min}\hfill
\perfil{Spike}{100}{130}{(0,0) (30,9) (40,100) (100,100) (130,0)}{0/0, 40/40, 100/100, 130/130 s}\hfill
\perfil{Stress}{400}{300}{(0,0) (60,100) (120,200) (180,300) (240,400) (300,0)}{0/0, 60/1, 120/2, 180/3, 240/4, 300/5 min}
\caption{Perfil de carga de cada escenario (usuarios virtuales en el tiempo).}
\end{figure}

%==================================================================
\section{Desarrollo de scripts en k6}
%==================================================================

\subsection{Herramientas y entorno}

\begin{table}[H]
\centering
\caption{Herramientas utilizadas.}
\begin{tabularx}{\textwidth}{l X}
\toprule
\rowcolor{ucbazul}
\encab{Herramienta} & \encab{Uso} \\
\midrule
Grafana k6 Studio 2.1.3 & Grabación del flujo y generación del script base (v1) \\
Grafana k6 (CLI)        & Ejecución de los scripts de carga \\
Visual Studio Code      & Edición del código \\
Git y GitHub            & Control de versiones y trabajo en equipo (ramas por paso) \\
Grafana k6 web dashboard & Visualización de métricas en tiempo real (\texttt{localhost:5665}) y reportes HTML \\
\bottomrule
\end{tabularx}
\end{table}

\subsection{Grabación del flujo con k6 Studio (v1)}

El recorrido se grabó con k6 Studio y se generó el script base con el asistente
de prueba. El script v1 contenía todas las peticiones del navegador con valores
fijos (tokens, cabeceras \texttt{traceparent}) y un solo check de status por petición.

\begin{pendiente}Capturas de la grabación y del script generado por el asistente.\end{pendiente}

\subsection{Correlación de valores dinámicos (v2 y v3)}

La correlación consiste en capturar valores que genera el servidor y reutilizarlos
en las peticiones siguientes. Se implementaron tres correlaciones:

\begin{table}[H]
\centering
\caption{Correlaciones implementadas.}
\begin{tabularx}{\textwidth}{c l l >{\raggedright\arraybackslash}X}
\toprule
\rowcolor{ucbazul}
\encab{\#} & \encab{Valor} & \encab{Técnica} & \encab{Origen y uso} \\
\midrule
1 & \texttt{csrf\_token} & Regex & Se extrae de la cabecera \texttt{Set-Cookie} de \texttt{POST /api/csrf-token} y se envía en el body del login \\
2 & Token de sesión & JSONPath (\texttt{token}) & Se extrae de la respuesta del login y se envía en la cabecera \texttt{Authorization} \\
3 & \texttt{pizza.id} & JSONPath (\texttt{pizza.id}) & Se extrae de \texttt{POST /api/pizza} y se envía en \texttt{POST /api/ratings} \\
\bottomrule
\end{tabularx}
\end{table}

\textbf{Hallazgos de la v1 a la v2:}
\begin{itemize}
  \item El asistente de k6 Studio detectó automáticamente la correlación del \texttt{csrf\_token}.
  \item No detectó la del \texttt{pizza\_id}: las calificaciones fallaban con \textbf{400}.
  \item Al correlacionarlo en modo JSON, k6 Studio lo enviaba como texto (\texttt{"20623"}) y el
        servidor (escrito en Go) esperaba un número entero. Se resolvió con un \emph{Replacer}
        en modo \emph{Begin-End}, obteniendo \textbf{201 Created}.
\end{itemize}

\textbf{Correcciones de la v2 a la v3:}
\begin{itemize}
  \item \emph{Love it} calificaba la pizza equivocada (usaba el ID de la primera). Ahora cada pizza guarda su propio ID.
  \item Los grupos \emph{No thanks} y \emph{Love it} repetían \texttt{POST /api/pizza} (ruido de la grabación),
        lo que inflaba el RPS. Se eliminaron.
  \item Se quitaron los tokens y \texttt{traceparent} fijos de la grabación.
  \item La extracción por regex ya no detiene la iteración si falta la cabecera.
\end{itemize}

Para evidenciar la correlación, el script tiene un \textbf{modo debug} que muestra
en la terminal los valores capturados en cada iteración:

\begin{gitbash}
(*\prompt*)
$ k6 run -e DEBUG=true --vus 1 --iterations 2 scripts/quickpizza-v3-modular.js
\end{gitbash}

\captura{images/correlacion-debug.png}{Valores correlacionados en dos iteraciones (modo debug).}

El \texttt{csrf\_token} y los \texttt{pizza.id} cambian en cada iteración. El token
de sesión del usuario \texttt{default} es fijo por diseño de QuickPizza, pero el
script no lo tiene escrito en el código: lo lee de la respuesta del login.

\subsection{Parametrización con SharedArray (v3)}
\label{sec:parametrizacion}

Los datos externos se cargan con \texttt{SharedArray}, que lee cada archivo
\textbf{una sola vez} y comparte los datos (solo lectura) entre todos los usuarios
virtuales. Con 100 VUs esto evita 100 copias de los mismos datos en memoria.

\begin{archivo}[language=JavaScript]{utils/data.js}
export const pizzas = new SharedArray("pizzas", function () {
  return parseCsv(open("../data/pizzas.csv")).map((p) => ({
    customName: p.customName,
    maxCaloriesPerSlice: Number(p.maxCaloriesPerSlice),
    mustBeVegetarian: p.mustBeVegetarian === "true",
    ...
  }));
});
\end{archivo}

\begin{itemize}
  \item \textbf{\texttt{data/pizzas.csv}:} 30 pizzas con nombre, calorías máximas por porción
        (700 a 1200), si es vegetariana, mínimo y máximo de toppings y herramienta excluida.
        Cada pedido usa una pizza al azar.
  \item \textbf{\texttt{data/users.csv}:} credenciales de inicio de sesión.
\end{itemize}

\textbf{Hallazgo sobre los usuarios:} se intentó usar 30 usuarios propios registrados
con \texttt{POST /api/users}. El registro respondía \emph{username already taken}
(el usuario existía), pero el login del \textbf{mismo} usuario a veces respondía 200
y a veces \textbf{401 authentication failed}. La hipótesis es que el sitio público
tiene varias réplicas y cada una guarda sus propios usuarios. Por eso se usa el
usuario \texttt{default}, que existe en todas las réplicas. El login se mantiene en
cada iteración como parte de la transacción completa.

\subsection{Validaciones de negocio con \texorpdfstring{\texttt{check()}}{check()}}

Cada iteración ejecuta 20 validaciones distintas, de tres tipos:

\begin{itemize}
  \item \textbf{Códigos de estado:} home, página de login, csrf, login, authenticate
        y pizza (200); calificación (201 Created).
  \item \textbf{Tiempos de respuesta:} la home, el csrf, el login, el pedido de pizza
        y la calificación deben responder en menos de 500\,ms
        (\texttt{r.timings.duration}). El límite está en \texttt{config/settings.js}
        y se puede cambiar con \texttt{-e MAX\_RESPONSE\_MS=...}.
  \item \textbf{Contenido de las respuestas:}
        \begin{itemize}
          \item la home devuelve la página HTML de QuickPizza;
          \item se extrae el \texttt{csrf\_token} y el token de sesión tiene 16 caracteres;
          \item el usuario autenticado es el correcto;
          \item la pizza recomendada tiene ID, nombre e ingredientes, y es vegetariana
                cuando el pedido lo exige (\texttt{mustBeVegetarian});
          \item la calificación guardada tiene el mismo \texttt{pizza\_id} y las mismas
                estrellas que se enviaron.
        \end{itemize}
\end{itemize}

Los checks del pedido de pizza y de la calificación se ejecutan dos veces por
iteración (una por pizza), por eso k6 reporta \textbf{28 checks por iteración}.

\begin{archivo}[language=JavaScript]{scenarios/orderPizza.js (checks de la calificación)}
check(res, {
  "rating: status 201": (r) => r.status === 201,
  [`rating: responde en < ${MAX_RESPONSE_MS} ms`]: respondeRapido,
  "rating: guarda la pizza y las estrellas": (r) =>
    r.status === 201 &&
    r.json("pizza_id") === pizzaId &&
    r.json("stars") === stars,
});
\end{archivo}

\textbf{Nota:} las corridas V5 de la sección~\ref{sec:resultados} se ejecutaron con
la versión anterior de los checks (14 por iteración: códigos de estado,
correlación y usuario autenticado), en la que se cumplió el 100\,\%. Los checks
de tiempo y de contenido se agregaron después.

\subsection{Definición de thresholds}

Los umbrales provisionales heredados de k6 Studio (\texttt{http\_req\_duration}
p(95) < 2100\,ms y p(99) < 2500\,ms) eran demasiado permisivos: los tiempos
reales rondan los 200\,ms. En la versión V5 se ajustaron a los criterios de
éxito, tanto en los cuatro escenarios nuevos como en el script v3:

\begin{archivo}[language=JavaScript]{scripts/quickpizza-v5-*.js (options)}
thresholds: {
  http_req_duration: ["p(95)<500"],
  http_req_failed: ["rate<0.01"],
},
\end{archivo}

Si se supera un umbral, k6 marca la prueba como fallida y termina con un código
de salida distinto de cero. Esto permite usar los scripts en un pipeline de CI.


\subsection{Estructura del repositorio}

El código está separado en datos, utilidades, configuración y escenarios:

\begin{archivo}{Estructura del repositorio}
quickpizza-performance-k6/
|-- config/settings.js        URL base, think time y modo debug
|-- data/users.csv            Credenciales
|-- data/pizzas.csv           30 pizzas
|-- utils/csv.js              Lector de CSV
|-- utils/data.js             SharedArray de usuarios y pizzas
|-- utils/http.js             Cabeceras y token anonimo
|-- utils/auth.js             csrf (regex), login y authenticate
|-- scenarios/orderPizza.js   Flujo de negocio (6 grupos)
|-- scripts/                  Scripts por version:
|     quickpizza-v3-modular.js      carga sostenida (v3)
|     quickpizza-v5-smoke.js        smoke
|     quickpizza-v5-load.js         carga
|     quickpizza-v5-spike.js        spike
|     quickpizza-v5-stress.js       estres
|-- reports/                  Reportes de k6
|-- docs/                     Informe y capturas
`-- README.md
\end{archivo}

El trabajo se organiza con una rama por paso (\texttt{feature/pasoXX-tema} o
\texttt{docs/tema}). Cada rama se revisa y se prueba antes de integrarla a
\texttt{main}. Una vez integrada, se borra tanto en local como en GitHub. Por
ejemplo, la rama \texttt{docs/update-readme-v5} se revisó con
\texttt{git show --stat}, sus cuatro escenarios se ejecutaron en local y luego se
integró a \texttt{main} con un \emph{fast-forward}, sin conflictos.

%==================================================================
\section{Ejecución y observabilidad con Grafana}
%==================================================================

\subsection{Configuración del envío de métricas}

Para monitorear las pruebas en tiempo real se utilizó el \textbf{dashboard web
integrado de Grafana k6}, que se ejecuta en el mismo equipo de pruebas
(\texttt{http://localhost:5665}). Esta opción cumple con trabajar en
\textbf{entorno local} sin instalar servicios adicionales ni usar Docker: k6 envía
sus métricas al dashboard cada 10 segundos mientras corre la prueba y, al terminar,
exporta un \textbf{reporte HTML interactivo} con las mismas gráficas en la carpeta
\texttt{reports/} del repositorio.

Cada escenario se ejecutó desde Git Bash con el dashboard habilitado. La salida de
consola se guardó además en un archivo de texto con \texttt{tee}, como evidencia:

\begin{gitbash}
(*\prompt*)
$ K6_WEB_DASHBOARD=true K6_WEB_DASHBOARD_EXPORT=reports/reporte-load-100vus.html \
  k6 run scripts/quickpizza-v5-load.js 2>&1 | tee reports/consola-load-100vus.txt

(*\prompt*)
$ K6_WEB_DASHBOARD=true K6_WEB_DASHBOARD_EXPORT=reports/reporte-spike.html \
  k6 run scripts/quickpizza-v5-spike.js 2>&1 | tee reports/consola-spike.txt

(*\prompt*)
$ K6_WEB_DASHBOARD=true K6_WEB_DASHBOARD_EXPORT=reports/reporte-stress-400vus.html \
  k6 run scripts/quickpizza-v5-stress.js
\end{gitbash}

Cada petición HTTP de \texttt{scenarios/\allowbreak orderPizza.js} y \texttt{utils/http.js}
lleva la etiqueta \texttt{tags.name} con el nombre del endpoint (por ejemplo,
\texttt{GET /}, \texttt{POST /api/csrf-token} o \texttt{POST /api/pizza}), lo que
permite agrupar las métricas por endpoint en el dashboard y en el resumen de k6.

La sintaxis \texttt{VARIABLE=valor comando} es de Bash (Git Bash). En PowerShell
las variables se definen antes con \texttt{\$env:K6\_WEB\_DASHBOARD="true"}.

El dashboard muestra las métricas clave que pide el análisis: tasa de peticiones
por segundo (RPS), distribución de latencias (promedio, P90, P95 y P99), tasa de
errores y número de usuarios virtuales activos, además del desglose de tiempos de
cada petición (conexión, TLS, envío, espera y recepción).

\subsection{Dashboard bajo carga máxima: RPS, latencias, errores y VUs}

\bloque{A}{Carga nominal sostenida -- Load (100 VUs)}

El escenario subió a 100 usuarios virtuales en el primer minuto y los
\textbf{sostuvo durante unos 3 minutos} (del minuto 1:00 al 4:02), para luego bajar.

\begin{itemize}
  \item \textbf{VUs y throughput:} con la meseta de 100 VUs el sistema procesó una
        tasa estable, con un promedio global de \textbf{68.8 req/s} y
        \textbf{21\,619 peticiones} completadas en 1\,663 iteraciones.
  \item \textbf{Latencias:} mediana de 218\,ms, P90 de 243\,ms, \textbf{P95 de
        248.5\,ms} y \textbf{P99 de 266.4\,ms}, muy por debajo del umbral de 500\,ms.
        La curva se mantuvo plana alrededor de 220\,ms durante toda la meseta.
  \item \textbf{Errores y validaciones:} 0\,\% de peticiones fallidas y 100\,\% de
        checks exitosos (23\,282 de 23\,282).
\end{itemize}

\captura{images/load-overview.png}{Dashboard Overview bajo carga nominal sostenida de 100 VUs (Request Rate, P95, VUs y Transfer Rate).}
\captura{images/load-vus-latencias.png}{Load: 100 VUs activos, duración de peticiones HTTP (P90/P95/P99) y duración de iteración.}
\captura{images/load-timings-1.png}{Load, pestaña Timings: Request Duration, Request Failed Rate, Request Rate y Request Waiting.}
\captura{images/load-timings-2.png}{Load, pestaña Timings: TLS Handshaking, Request Sending, Connecting y Receiving.}
\captura{images/load-summary.png}{Load: resumen estadístico (Summary) del reporte HTML.}

\bloque{B}{Pico de usuarios -- Spike (100 VUs en 10 segundos)}

Tras 30 segundos de calentamiento con menos de 10 VUs, la carga saltó de
\textbf{9 a 100 VUs en 10 segundos}, se sostuvo durante un minuto y luego bajó.

\begin{itemize}
  \item \textbf{VUs y throughput:} 7\,800 peticiones en 600 iteraciones, con un
        promedio de 55\,req/s (la prueba es más corta y parte de ella transcurre con
        pocos usuarios).
  \item \textbf{Latencias:} mediana de 220\,ms, P90 de 245\,ms, \textbf{P95 de
        251.1\,ms} y \textbf{P99 de 290.2\,ms}. El sistema absorbió la llegada brusca de usuarios sin
        degradación visible.
  \item \textbf{Errores y validaciones:} 0\,\% de errores y 100\,\% de checks
        exitosos (8\,400 de 8\,400).
\end{itemize}

\captura{images/spike-overview.png}{Dashboard Overview del escenario spike: salto de 9 a 100 VUs en 10 segundos.}
\captura{images/spike-vus-latencias.png}{Spike: forma del pico de VUs y comportamiento de las latencias.}
\captura{images/spike-timings-1.png}{Spike, pestaña Timings: Request Duration, Request Failed Rate, Request Rate y Request Waiting.}
\captura{images/spike-summary.png}{Spike: resumen estadístico (Summary) del reporte HTML.}

\bloque{C}{Carga máxima de estrés -- Stress (400 VUs)}

La carga subió en escalones de un minuto (100, 200, 300 y 400 VUs) para evaluar
el comportamiento al \textbf{cuadruplicar la carga nominal}. Las capturas se
tomaron con el dashboard en vivo durante el tramo de máxima concurrencia.

\begin{itemize}
  \item \textbf{VUs y pico de RPS:} el throughput creció de forma pareja con los
        usuarios, hasta un pico de \textbf{unos 330 req/s} con 400 VUs (alrededor
        de las 20:44:50). En total se completaron unas 54\,500 peticiones y 4\,200
        iteraciones, con un promedio de 172\,req/s.
  \item \textbf{Latencias:} en el resultado global, mediana de 222\,ms, P90 de
        257\,ms, \textbf{P95 de 278\,ms} y \textbf{P99 de 374\,ms}. En el minuto de
        máxima concurrencia, el P99 subió hasta unos 470\,ms y el P95 hasta unos
        350\,ms; al bajar los usuarios, ambos volvieron a unos 250\,ms.
  \item \textbf{Origen del tiempo de respuesta:} \texttt{http\_req\_waiting}
        (P95 = 277\,ms) es prácticamente igual a \texttt{http\_req\_duration}
        (P95 = 278\,ms), mientras que \texttt{http\_req\_receiving} es mínimo
        (P95 = 1\,ms). Casi todo el tiempo corresponde al \textbf{procesamiento en el
        servidor}, no a la red ni a la descarga de datos.
  \item \textbf{Errores y validaciones:} 0\,\% de errores y 100\,\% de checks
        exitosos.
\end{itemize}

\captura{images/stress-overview.png}{Dashboard Overview bajo estrés: el RPS sube con los VUs hasta unos 310--330 req/s y 0\,\% de errores.}
\captura{images/stress-vus-latencias.png}{Stress: escalera hasta 400 VUs y su bajada, con el aumento de P95/P99 en el minuto de máxima concurrencia.}
\captura{images/stress-timings-1.png}{Stress, pestaña Timings: Request Duration y Request Waiting siguen la misma curva (tiempo de procesamiento en el servidor).}
\captura{images/stress-timings-2.png}{Stress, pestaña Timings: TLS y Connecting solo al inicio (conexiones reutilizadas); Receiving se mantiene en pocos milisegundos.}
\captura{images/stress-summary.png}{Stress: resumen estadístico (Summary): 54.5k peticiones, P90 = 257\,ms, P95 = 278\,ms, P99 = 374\,ms y 0\,\% de errores.}

\subsection{Evidencias de ejecución}
\label{sec:resultados}

Los escenarios se ejecutaron el 18 de noviembre de 2026 contra \url{\appURL}.
Todos pasaron los umbrales:

\begin{table}[H]
\centering
\caption{Resultados de los escenarios V5.}
\begin{tabularx}{\textwidth}{X r r r r r r r r}
\toprule
\rowcolor{ucbazul}
\encab{Escenario} & \encab{VUs} & \encab{Peticiones} & \encab{Req/s} & \encab{Mediana} & \encab{P95} & \encab{P99} & \encab{Máx.} & \encab{Errores} \\
\midrule
\esmoke  & 1   & 13      & 0.8   & 205\,ms & 225\,ms   & --       & 226\,ms   & 0\,\% \\
\eload   & 100 & 21\,619 & 68.8  & 218\,ms & 248.5\,ms & 266.4\,ms & 1\,240\,ms & 0\,\% \\
\espike  & 100 & 7\,800  & 55.0  & 220\,ms & 251.1\,ms & 290.2\,ms & 1\,125\,ms & 0\,\% \\
\estress & 400 & 54\,500 & 172.0 & 222\,ms & 278\,ms   & 374\,ms  & 1\,000\,ms & 0\,\% \\
\bottomrule
\end{tabularx}
\end{table}


\begin{table}[H]
\centering
\caption{Iteraciones, checks y calificaciones por escenario.}
\begin{tabularx}{\textwidth}{l >{\hspace{1.2cm}}r >{\centering\arraybackslash}X r >{\hspace{0.6cm}}r}
\toprule
\rowcolor{ucbazul}
\encab{Escenario} & \encab{Iteraciones} & \encab{Checks OK} & \encab{Love it} & \encab{No thanks} \\
\midrule
\esmoke  & 1       & 14 / 14           & 2       & 0       \\
\eload   & 1\,663  & 23\,282 / 23\,282 & 2\,678  & 648     \\
\espike  & 600     & 8\,400 / 8\,400   & 964     & 236     \\
\estress & 4\,200  & 100\,\%           & 6\,800  & 1\,600  \\
\bottomrule
\end{tabularx}
\end{table}

En todos los escenarios, \texttt{love\_it} + \texttt{no\_thanks} es el doble de las
iteraciones: ninguna iteración se interrumpió y todas las calificaciones quedaron
registradas. Los valores del stress están redondeados por el dashboard.

\textbf{Repetibilidad del escenario de estrés.} El stress de 400 VUs se ejecutó tres
veces durante el proyecto, desde distintos equipos y en distintos momentos:

\begin{table}[H]
\centering
\caption{Ejecuciones del escenario stress (400 VUs).}
\begin{tabularx}{\textwidth}{X r r l}
\toprule
\rowcolor{ucbazul}
\encab{Ejecución} & \encab{P95} & \encab{Errores} & \encab{Resultado} \\
\midrule
Corrida 1 (mañana)                   & 270\,ms & 0\,\% & \aprobado \\
Corrida 2 (otro integrante del equipo) & 584\,ms & 0\,\% & \noaprobado{(P95 > 500\,ms)} \\
Corrida 3 (noche, con dashboard)      & 278\,ms & 0\,\% & \aprobado \\
\bottomrule
\end{tabularx}
\end{table}

En ninguna de las tres hubo errores HTTP, pero en una el P95 superó el umbral. Con
400 VUs el sistema trabaja cerca de su límite, y el resultado depende de las
condiciones del momento (QuickPizza es un servicio público compartido con otros
usuarios) y de la red desde donde se genera la carga. Las capturas de este informe
corresponden a la corrida 3.

\captura{images/v5-smoke.png}{Salida de consola del escenario smoke.}
\captura{images/load-consola.png}{Salida de consola del escenario load (100 VUs).}
\captura{images/spike-consola.png}{Salida de consola del escenario spike (100 VUs).}

%==================================================================
\section{Análisis de resultados}
%==================================================================

\subsection{Interpretación de gráficas}

\begin{itemize}
  \item \textbf{Smoke:} el flujo completo funciona con un usuario, con respuestas
        de 200 a 226\,ms. Sirve como línea base.
  \item \textbf{Load (100 VUs):} pasar de 1 a 100 usuarios casi no cambió la latencia
        (P95 de 225 a 248.5\,ms) y la curva se mantuvo plana durante los 3 minutos de
        meseta. La carga nominal se maneja con amplio margen frente al umbral.
  \item \textbf{Spike (100 VUs en 10 s):} la llegada brusca de usuarios no causó
        errores ni degradación: el P95 (251.1\,ms) es prácticamente igual al del
        load. El sistema absorbe picos repentinos de esta magnitud.
  \item \textbf{Stress (400 VUs):} el throughput escala casi en proporción a los
        usuarios (hasta unos 330 req/s), pero en el escalón de 400 VUs aparece una
        \textbf{degradación leve y reversible}: el P99 sube a unos 470\,ms y el P95 a
        unos 350\,ms durante el pico, y vuelven a unos 250\,ms cuando bajan los
        usuarios. La mediana casi no cambia (222\,ms), así que la mayoría de los
        usuarios no lo nota.
\end{itemize}

En la pestaña Timings del stress, la conexión TCP y el \emph{handshake} TLS solo
tienen picos al inicio, cuando entran usuarios nuevos; después quedan en 0\,ms
porque k6 reutiliza las conexiones (\emph{keep-alive}).

\subsection{Cuellos de botella identificados}

\begin{hallazgo}
{\sffamily\bfseries\color{ucbazulclaro} Hallazgo clave.} Casi todo el tiempo de respuesta
corresponde al procesamiento en el servidor, y las operaciones que escriben datos
(pedir y calificar pizzas) son las más lentas. Alrededor de 300--400 usuarios
concurrentes las peticiones empiezan a hacer cola, pero sin errores.
\end{hallazgo}

Como QuickPizza es un servicio público, no se tiene acceso a sus métricas internas
(CPU, memoria, base de datos). El diagnóstico se basa en lo que se observa desde el
cliente:

\begin{itemize}
  \item \textbf{Procesamiento en el servidor.} En todos los escenarios el tiempo de
        espera (\texttt{http\_req\_waiting}) explica casi todo el tiempo de respuesta,
        y el de recepción es de pocos milisegundos. El cuello de botella está en el
        procesamiento del backend (lógica de la aplicación o acceso a datos), no en
        la red.
  \item \textbf{Las escrituras son las operaciones más lentas.} El reporte HTML del
        load desglosa los tiempos por grupo del flujo (tabla~\ref{tab:grupos}). Las
        operaciones que escriben en el servidor (\texttt{POST /api/pizza} y, sobre
        todo, \texttt{POST /api/ratings}) tardan unos 25--35\,ms más que las de
        lectura de la página de inicio y el login. Calificar una pizza, que guarda
        datos, es la operación más lenta (P99 de 285\,ms), lo que apunta al acceso a
        la base de datos como el componente más costoso del flujo.
  \item \textbf{Cercanía al límite con 400 VUs.} El aumento del P99 en el escalón
        de máxima concurrencia, y la corrida en la que el P95 llegó a 584\,ms, indican
        que alrededor de \textbf{300--400 usuarios concurrentes} las peticiones
        empiezan a hacer cola en el servidor. No se rechazó ninguna petición: es una
        degradación gradual, no una caída.
  \item \textbf{Carga efectiva limitada por el \emph{think time}.} Cada iteración
        dura unos 15\,s, casi todo de pausas, así que 100 VUs generan solo unos
        70\,req/s. El número de VUs sobreestima la presión real sobre el servidor.
  \item \textbf{Estado no compartido entre réplicas.} Los usuarios registrados no se
        replican entre las instancias del sitio, lo que provoca respuestas 401
        intermitentes en el login (ver sección~\ref{sec:parametrizacion}).
\end{itemize}

\begin{table}[H]
\centering
\caption{Duración de las peticiones por grupo del flujo (load, 100 VUs).}
\label{tab:grupos}
\begin{tabularx}{\textwidth}{X l r r r r}
\toprule
\rowcolor{ucbazul}
\encab{Grupo} & \encab{Tipo} & \encab{Peticiones} & \encab{Mediana} & \encab{P95} & \encab{P99} \\
\midrule
01. HomePage          & Lectura   & 6\,652 & 214\,ms & 227\,ms & 242\,ms \\
02. Login             & Lectura / sesión & 8\,315 & 215\,ms & 228\,ms & 247\,ms \\
03. Pedir pizza 1     & Escritura & 1\,663 & 234\,ms & 253\,ms & 274\,ms \\
04. Calificar pizza 1 & Escritura & 1\,663 & 242\,ms & 261\,ms & 284\,ms \\
05. Pedir pizza 2     & Escritura & 1\,663 & 235\,ms & 254\,ms & 279\,ms \\
06. Calificar pizza 2 & Escritura & 1\,663 & 242\,ms & 260\,ms & 285\,ms \\
\bottomrule
\end{tabularx}
\end{table}

\subsection{Recomendaciones de optimización}

\textbf{Para la aplicación:}
\begin{itemize}
  \item \textbf{Escalado horizontal:} añadir réplicas o configurar autoescalado
        cuando la concurrencia se acerque a unos 300 usuarios, el punto en el que
        empieza la degradación.
  \item \textbf{Caché:} las peticiones de la página de inicio
        (\texttt{/api/config}, \texttt{/api/quotes}, \texttt{/api/tools}) devuelven
        datos que casi no cambian; servirlas desde caché reduce la carga del backend.
  \item \textbf{Optimizar las escrituras en base de datos:} revisar las consultas de
        \texttt{/api/ratings} y \texttt{/api/pizza} y crear índices sobre los campos
        que se consultan al guardar una calificación (por ejemplo, el ID de la pizza
        y del usuario), ya que son las operaciones más lentas del flujo.
  \item \textbf{Sesiones compartidas:} guardar usuarios y sesiones en un
        almacenamiento común a todas las réplicas (base de datos o caché central), o
        usar \emph{sticky sessions} en el balanceador.
  \item \textbf{Monitoreo de la cola de latencias:} vigilar el P99 en producción,
        porque es donde aparecen los primeros síntomas de saturación.
\end{itemize}

\textbf{Para las pruebas:}
\begin{itemize}
  \item Subir el pico del \emph{spike} a 300--400 VUs para probar un pico mayor que
        la carga nominal.
  \item Agregar umbrales por endpoint (por ejemplo, el login) con etiquetas de k6,
        para detectar degradaciones localizadas.
  \item Ejecutar las pruebas de estrés pocas veces y de forma controlada, porque el
        servicio es compartido y los resultados varían según el momento.
\end{itemize}

%==================================================================
\section{Conclusiones}
%==================================================================

\begin{itemize}
  \item QuickPizza cumple los criterios de éxito bajo carga nominal: con 100 VUs el
        P95 fue de 248.5\,ms (umbral de 500\,ms), con 0\,\% de errores y 100\,\% de
        checks. También superó el pico repentino de 100 VUs en 10 segundos.
  \item Con 400 VUs (cuatro veces la carga nominal) el sistema sigue estable y sin
        errores, pero se acerca a su límite: el P99 sube durante el pico y, en una de
        tres ejecuciones, el P95 superó los 500\,ms. La capacidad recomendada se ubica
        en torno a los 300 usuarios concurrentes.
  \item El tiempo de respuesta depende casi por completo del procesamiento en el
        servidor; la red no es un factor limitante.
  \item La correlación (CSRF, token de sesión y \texttt{pizza.id}) y la
        parametrización con \texttt{SharedArray} funcionan correctamente bajo carga:
        miles de iteraciones completas sin un solo fallo de negocio.
  \item La estructura modular permitió crear los escenarios V5 reutilizando el mismo
        flujo (\texttt{scenarios/\allowbreak orderPizza.js}), cambiando solo las \texttt{options}.
\end{itemize}

\begin{pendiente}Revisar y completar las conclusiones con el equipo.\end{pendiente}

%==================================================================
\section{Anexos}
%==================================================================

\subsection{Repositorio y ejecución de las pruebas}

Repositorio: \url{\repoURL}

\begin{gitbash}
(*\prompt*)
$ # Prueba rapida (1 usuario)
$ k6 run --vus 1 --iterations 1 scripts/quickpizza-v3-modular.js

$ # Ver las correlaciones
$ k6 run -e DEBUG=true --vus 1 --iterations 2 scripts/quickpizza-v3-modular.js

$ # Carga completa (100 VUs)
$ k6 run scripts/quickpizza-v3-modular.js

$ # Escenarios V5
$ k6 run scripts/quickpizza-v5-smoke.js
$ k6 run scripts/quickpizza-v5-load.js
$ k6 run scripts/quickpizza-v5-spike.js
$ k6 run scripts/quickpizza-v5-stress.js
\end{gitbash}

\subsection{Reportes de ejecución en el repositorio}

La evidencia de las corridas está en la carpeta \texttt{reports/} del repositorio:

\begin{itemize}
  \item \texttt{reporte-load-100vus.html}, \texttt{reporte-spike.html} y
        \texttt{reporte-stress-400vus.html}: reportes interactivos del dashboard de
        Grafana k6, con las mismas gráficas de este informe. Se abren descargándolos
        y abriéndolos en un navegador.
  \item \texttt{consola-load-100vus.txt} y \texttt{consola-spike.txt}: salida
        completa de la consola de k6, con el resumen final de cada prueba.
\end{itemize}

\subsection{Integrantes y acceso al repositorio}

El repositorio es público y el propietario agregó como colaboradores a los
integrantes del equipo:

\begin{table}[H]
\centering
\caption{Integrantes del equipo y sus cuentas de GitHub.}
\begin{tabularx}{\textwidth}{X l l}
\toprule
\rowcolor{ucbazul}
\encab{Integrante} & \encab{Usuario de GitHub} & \encab{Rol en el repositorio} \\
\midrule
Pablo Ivan Chuquimia Huanca    & \texttt{pchuquimia}      & Colaborador \\
Mauricio Gutierrez Pacheco     & \texttt{ElNegis}         & Colaborador \\
Rodrigo Camilo Mendez Gallardo & \texttt{CamiloMendez23}  & Colaborador \\
Maria Alexandra Quintanilla    & \texttt{Marialejandra15} & Colaboradora \\
Pedro Adrian Santos Orellana   & \texttt{nohpedro}        & Colaborador \\
David Soliz Garcia             & \texttt{Soliz-G-David}   & Propietario \\
\bottomrule
\end{tabularx}
\end{table}

\subsection{Historial de commits}

Historial de la rama \texttt{main} (obtenido con \texttt{git log}). Cada paso se
trabajó en una rama propia, se revisó y se integró a \texttt{main}:

\begin{table}[H]
\centering
\caption{Commits del proyecto en la rama \texttt{main}.}
\begin{tabularx}{\textwidth}{l l l >{\raggedright\arraybackslash}X}
\toprule
\rowcolor{ucbazul}
\encab{Commit} & \encab{Fecha} & \encab{Autor} & \encab{Descripción} \\
\midrule
\texttt{a66bc6e} & 16/11/2026 & David Soliz & Estructura base del proyecto + script v2 generado con k6 Studio (correlaciones csrf y \texttt{pizza\_id}) \\
\texttt{e948877} & 17/11/2026 & David Soliz & [Paso 1] 2.c Parametrización: archivos CSV con credenciales y 30 pizzas \\
\texttt{e8284a8} & 17/11/2026 & David Soliz & [Paso 1] 2.c Parametrización: carga de CSV con \texttt{SharedArray} para ahorrar memoria en los VUs \\
\texttt{58cdcc9} & 17/11/2026 & David Soliz & [Paso 2] 2.c Correlación: csrf (regex), token de sesión y \texttt{pizza.id} (JSONPath) + script v3 modular \\
\texttt{dfd5b63} & 17/11/2026 & Alexandra Quintanilla & Distribución de calificaciones (v4) \\
\texttt{56ce305} & 17/11/2026 & Alexandra Quintanilla & Escenarios V5 (smoke, load, spike, stress), umbrales p(95) < 500\,ms y README actualizado \\
\texttt{3cba42c} & 18/11/2026 & David Soliz & Reportes de ejecución (salida de consola de load y spike) \\
\texttt{7c7286b} & 18/11/2026 & David Soliz & Reportes HTML del dashboard k6 (load, spike y stress) \\
\bottomrule
\end{tabularx}
\end{table}

\subsection{Scripts de k6}

Código final de los archivos principales. El repositorio completo está en
\url{\repoURL}.

\begin{archivo}[language=JavaScript]{scenarios/orderPizza.js}
// =====================================================================
// scenarios/orderPizza.js
// Flujo de negocio: un usuario entra a QuickPizza, inicia sesión,
// pide una pizza, la califica, pide otra y la vuelve a calificar.
//
// V4:
//  - Parametrización mediante CSV.
//  - Conversión de tipos correctos para el backend.
//  - Correlación de CSRF, sesión y pizza.id.
//  - Cada pizza conserva su propio ID.
//  - Distribución de calificaciones: 80% Love it / 20% No thanks.
//  - Contadores para medir la distribución de Love it / No thanks.
//
// V6:
//  - Checks de tiempo de respuesta (< MAX_RESPONSE_MS) en cada petición validada.
//  - Checks de contenido: HTML de la home, nombre e ingredientes de la pizza,
//    pizza vegetariana cuando se pide y datos de la calificación registrada.
// =====================================================================

import http from "k6/http";
import { group, check, sleep } from "k6";
import { Counter } from "k6/metrics";

import { BASE_URL, THINK_TIME, MAX_RESPONSE_MS, debugLog } from "../config/settings.js";
import { htmlHeaders, apiParams, randomToken } from "../utils/http.js";
import { getCsrfToken, login, authenticate } from "../utils/auth.js";
import { getUserForVU, getRandomPizza } from "../utils/data.js";

// Contadores para verificar la distribución 80/20.
const loveItCounter = new Counter("love_it");
const noThanksCounter = new Counter("no_thanks");

// Check de tiempo: la petición respondió antes del límite configurado.
function respondeRapido(r) {
  return r.timings.duration < MAX_RESPONSE_MS;
}

function thinkTime() {
  sleep(
    THINK_TIME.min +
      Math.random() * (THINK_TIME.max - THINK_TIME.min)
  );
}

// Determina aleatoriamente si la pizza recibe Love it.
// Aproximadamente 80% = Love it
// Aproximadamente 20% = No thanks
function shouldLoveIt() {
  return Math.random() < 0.8;
}

// Pide una pizza con los datos del CSV y devuelve su ID.
// Correlación JSONPath: pizza.id
function orderPizza(token, pizzaData) {
  const res = http.post(
    `${BASE_URL}/api/pizza`,
    JSON.stringify(pizzaData),
    apiParams(token, "POST /api/pizza"),
  );

  const pizzaId = res.status === 200 ? res.json("pizza.id") : null;

  debugLog(
    `Correlación 3 (JSONPath) pizza.id ("${pizzaData.customName}")`,
    pizzaId
  );

  // Lee el body una sola vez; si no es JSON válido, queda en null.
  let body = null;
  try {
    body = res.json();
  } catch (e) {
    body = null;
  }
  const pizza = body && body.pizza ? body.pizza : null;

  check(res, {
    "pizza: status 200": (r) => r.status === 200,
    "pizza: tiene id": () => pizzaId !== null,
    [`pizza: responde en < ${MAX_RESPONSE_MS} ms`]: respondeRapido,
    "pizza: tiene nombre e ingredientes": () =>
      pizza !== null &&
      typeof pizza.name === "string" &&
      pizza.name.length > 0 &&
      Array.isArray(pizza.ingredients) &&
      pizza.ingredients.length > 0,
    "pizza: es vegetariana si se pidió": () =>
      !pizzaData.mustBeVegetarian || (body !== null && body.vegetarian === true),
  });

  return pizzaId;
}

// Califica una pizza.
// pizza_id debe enviarse como número porque el servidor Go espera int64.
function ratePizza(token, pizzaId, stars) {
  const res = http.post(
    `${BASE_URL}/api/ratings`,
    JSON.stringify({
      pizza_id: pizzaId,
      stars: stars,
    }),
    apiParams(token, "POST /api/ratings"),
  );

  check(res, {
    "rating: status 201": (r) => r.status === 201,
    [`rating: responde en < ${MAX_RESPONSE_MS} ms`]: respondeRapido,
    "rating: guarda la pizza y las estrellas": (r) =>
      r.status === 201 &&
      r.json("pizza_id") === pizzaId &&
      r.json("stars") === stars,
  });

  return res;
}

export function orderPizzaFlow() {
  const user = getUserForVU();
  const anonToken = randomToken();
  let token = null;

  group("01. HomePage", function () {
    const res = http.get(
      `${BASE_URL}/`,
      Object.assign(htmlHeaders(), {
        tags: { name: "GET /" },
      }),
    );

    check(res, {
      "home: status 200": (r) => r.status === 200,
      [`home: responde en < ${MAX_RESPONSE_MS} ms`]: respondeRapido,
      "home: devuelve la página de QuickPizza": (r) =>
        typeof r.body === "string" && r.body.includes("<html") && r.body.includes("pizza"),
    });

    http.get(
      `${BASE_URL}/api/config`,
      apiParams(null, "GET /api/config")
    );

    http.get(
      `${BASE_URL}/api/quotes`,
      apiParams(null, "GET /api/quotes")
    );

    http.get(
      `${BASE_URL}/api/tools`,
      apiParams(anonToken, "GET /api/tools")
    );
  });

  thinkTime();

  group("02. Login", function () {
    const res = http.get(
      `${BASE_URL}/login`,
      Object.assign(htmlHeaders(`${BASE_URL}/`), {
        tags: { name: "GET /login" },
      }),
    );

    check(res, {
      "login page: status 200": (r) => r.status === 200,
    });

    const csrfToken = getCsrfToken();

    // Correlación 1: Regex
    if (!csrfToken) return;

    // Correlación 2: JSONPath
    token = login(user, csrfToken);

    if (!token) return;

    authenticate(token, user.username);

    http.get(
      `${BASE_URL}/api/ratings`,
      apiParams(token, "GET /api/ratings")
    );
  });

  // Si el login falla, no tiene sentido continuar.
  if (!token) return;

  thinkTime();

  let firstPizzaId = null;

  group("03. Pedir pizza 1", function () {
    firstPizzaId = orderPizza(
      token,
      getRandomPizza()
    );
  });

  thinkTime();

  group("04. Calificar pizza 1", function () {
    if (firstPizzaId !== null) {
      const loveIt = shouldLoveIt();

      if (loveIt) {
        loveItCounter.add(1);
        ratePizza(token, firstPizzaId, 5);
      } else {
        noThanksCounter.add(1);
        ratePizza(token, firstPizzaId, 1);
      }
    }
  });

  thinkTime();

  let secondPizzaId = null;

  group("05. Pedir pizza 2", function () {
    secondPizzaId = orderPizza(
      token,
      getRandomPizza()
    );
  });

  thinkTime();

  group("06. Calificar pizza 2", function () {
    if (secondPizzaId !== null) {
      const loveIt = shouldLoveIt();

      if (loveIt) {
        loveItCounter.add(1);
        ratePizza(token, secondPizzaId, 5);
      } else {
        noThanksCounter.add(1);
        ratePizza(token, secondPizzaId, 1);
      }
    }
  });

  thinkTime();
}
\end{archivo}

\begin{archivo}[language=JavaScript]{utils/auth.js}
// =====================================================================
// utils/auth.js
// Autenticación y correlación de tokens.
//
// Correlación 1 (Regex):    csrf_token  -> se extrae de la cabecera Set-Cookie
// Correlación 2 (JSONPath): token       -> se extrae del body JSON del login
// =====================================================================

import http from "k6/http";
import { check } from "k6";
import { BASE_URL, MAX_RESPONSE_MS, debugLog } from "../config/settings.js";
import { apiParams } from "./http.js";

// ---------------------------------------------------------------------
// POST /api/csrf-token
// El servidor responde con: Set-Cookie: csrf_token=XXXX; Path=/; SameSite=Strict
// ---------------------------------------------------------------------
export function getCsrfToken() {
  const res = http.post(`${BASE_URL}/api/csrf-token`, null, {
    headers: { accept: "*/*", origin: BASE_URL, referer: `${BASE_URL}/login` },
    tags: { name: "POST /api/csrf-token" },
  });

  // Se valida que la cabecera exista antes de aplicar la regex
  // (en la v2 un header vacío rompía la iteración).
  const setCookie = res.headers["Set-Cookie"] || "";
  const match = setCookie.match(/csrf_token=([A-Za-z0-9]+);/);
  const csrfToken = match ? match[1] : null;
  debugLog("Correlación 1 (Regex)    csrf_token", csrfToken);

  check(res, {
    "csrf: status 200": (r) => r.status === 200,
    "csrf: token extraído": () => csrfToken !== null,
    [`csrf: responde en < ${MAX_RESPONSE_MS} ms`]: (r) => r.timings.duration < MAX_RESPONSE_MS,
  });

  return csrfToken;
}

// ---------------------------------------------------------------------
// POST /api/users/token/login?set_cookie=true
// El csrf enviado en el body debe coincidir con la cookie csrf_token.
// Respuesta: { "token": "XXXXXXXXXXXXXXXX" }  -> se extrae con JSONPath
// ---------------------------------------------------------------------
export function login(user, csrfToken) {
  const params = apiParams(null, "POST /api/users/token/login", `${BASE_URL}/login`);
  // Se envía la cookie csrf_token de forma explícita (igual que en la v2):
  // el servidor compara esta cookie con el campo "csrf" del body.
  params.cookies = {
    csrf_token: { value: csrfToken, replace: true },
  };

  const res = http.post(
    `${BASE_URL}/api/users/token/login?set_cookie=true`,
    JSON.stringify({
      username: user.username,
      password: user.password,
      csrf: csrfToken,
    }),
    params,
  );

  const token = res.status === 200 ? res.json("token") : null;
  debugLog("Correlación 2 (JSONPath) token de sesión", token);

  if (res.status !== 200) {
    console.warn(`[login] ${user.username} -> status ${res.status} | body: ${res.body}`);
  }

  check(res, {
    "login: status 200": (r) => r.status === 200,
    "login: token de 16 caracteres": () => token !== null && token.length === 16,
    [`login: responde en < ${MAX_RESPONSE_MS} ms`]: (r) => r.timings.duration < MAX_RESPONSE_MS,
  });

  return token;
}

// ---------------------------------------------------------------------
// POST /api/users/token/authenticate
// Confirma que el token capturado es válido y pertenece al usuario.
// ---------------------------------------------------------------------
export function authenticate(token, username) {
  const res = http.post(
    `${BASE_URL}/api/users/token/authenticate`,
    null,
    apiParams(token, "POST /api/users/token/authenticate"),
  );

  check(res, {
    "authenticate: status 200": (r) => r.status === 200,
    "authenticate: usuario correcto": (r) => r.status === 200 && r.json("username") === username,
  });

  return res;
}
\end{archivo}

\begin{archivo}[language=JavaScript]{config/settings.js}
// =====================================================================
// config/settings.js
// Configuración general del proyecto (URL base, tiempos de espera, debug).
// La URL se puede cambiar sin tocar el código:
//   k6 run -e BASE_URL=http://localhost:3333 scripts/quickpizza-v3-modular.js
// =====================================================================

export const BASE_URL = __ENV.BASE_URL || "https://quickpizza.grafana.com";

// Think time (segundos) entre pasos del usuario: simula a una persona real.
export const THINK_TIME = {
  min: 1,
  max: 3,
};

// Modo debug: muestra en la terminal los valores correlacionados.
//   k6 run -e DEBUG=true --vus 1 --iterations 2 scripts/quickpizza-v3-modular.js
export const DEBUG = __ENV.DEBUG === "true";

// Tiempo máximo de respuesta (ms) que validan los checks de cada petición.
// Coincide con el criterio de éxito (P95 < 500 ms). Se puede cambiar sin tocar el código:
//   k6 run -e MAX_RESPONSE_MS=800 scripts/quickpizza-v5-load.js
export const MAX_RESPONSE_MS = Number(__ENV.MAX_RESPONSE_MS) || 500;

export function debugLog(label, value) {
  if (DEBUG) {
    console.log(`[VU ${__VU} | iteración ${__ITER}] ${label} = ${value}`);
  }
}
\end{archivo}

\begin{archivo}[language=JavaScript]{scripts/quickpizza-v5-stress.js}
import { orderPizzaFlow } from "../scenarios/orderPizza.js";

export const options = {
  stages: [
    { target: 100, duration: "1m" },
    { target: 200, duration: "1m" },
    { target: 300, duration: "1m" },
    { target: 400, duration: "1m" },
    { target: 0, duration: "1m" },
  ],

  thresholds: {
    http_req_duration: ["p(95)<500"],
    http_req_failed: ["rate<0.01"],
  },
};

export default function () {
  orderPizzaFlow();
}
\end{archivo}

\end{document}